% Options for packages loaded elsewhere
% Options for packages loaded elsewhere
\PassOptionsToPackage{unicode}{hyperref}
\PassOptionsToPackage{hyphens}{url}
\PassOptionsToPackage{dvipsnames,svgnames,x11names}{xcolor}
%
\documentclass[
  canadian,
]{article}
\usepackage{xcolor}
\usepackage[top = 0.75in,bottom = 1in,left = 1in,right = 1in]{geometry}
\usepackage{amsmath,amssymb}
\setcounter{secnumdepth}{-\maxdimen} % remove section numbering
\usepackage{iftex}
\ifPDFTeX
  \usepackage[T1]{fontenc}
  \usepackage[utf8]{inputenc}
  \usepackage{textcomp} % provide euro and other symbols
\else % if luatex or xetex
  \usepackage{unicode-math} % this also loads fontspec
  \defaultfontfeatures{Scale=MatchLowercase}
  \defaultfontfeatures[\rmfamily]{Ligatures=TeX,Scale=1}
\fi
\usepackage{lmodern}
\ifPDFTeX\else
  % xetex/luatex font selection
\fi
% Use upquote if available, for straight quotes in verbatim environments
\IfFileExists{upquote.sty}{\usepackage{upquote}}{}
\IfFileExists{microtype.sty}{% use microtype if available
  \usepackage[]{microtype}
  \UseMicrotypeSet[protrusion]{basicmath} % disable protrusion for tt fonts
}{}
\makeatletter
\@ifundefined{KOMAClassName}{% if non-KOMA class
  \IfFileExists{parskip.sty}{%
    \usepackage{parskip}
  }{% else
    \setlength{\parindent}{0pt}
    \setlength{\parskip}{6pt plus 2pt minus 1pt}}
}{% if KOMA class
  \KOMAoptions{parskip=half}}
\makeatother
% Make \paragraph and \subparagraph free-standing
\makeatletter
\ifx\paragraph\undefined\else
  \let\oldparagraph\paragraph
  \renewcommand{\paragraph}{
    \@ifstar
      \xxxParagraphStar
      \xxxParagraphNoStar
  }
  \newcommand{\xxxParagraphStar}[1]{\oldparagraph*{#1}\mbox{}}
  \newcommand{\xxxParagraphNoStar}[1]{\oldparagraph{#1}\mbox{}}
\fi
\ifx\subparagraph\undefined\else
  \let\oldsubparagraph\subparagraph
  \renewcommand{\subparagraph}{
    \@ifstar
      \xxxSubParagraphStar
      \xxxSubParagraphNoStar
  }
  \newcommand{\xxxSubParagraphStar}[1]{\oldsubparagraph*{#1}\mbox{}}
  \newcommand{\xxxSubParagraphNoStar}[1]{\oldsubparagraph{#1}\mbox{}}
\fi
\makeatother


\usepackage{longtable,booktabs,array}
\usepackage{calc} % for calculating minipage widths
% Correct order of tables after \paragraph or \subparagraph
\usepackage{etoolbox}
\makeatletter
\patchcmd\longtable{\par}{\if@noskipsec\mbox{}\fi\par}{}{}
\makeatother
% Allow footnotes in longtable head/foot
\IfFileExists{footnotehyper.sty}{\usepackage{footnotehyper}}{\usepackage{footnote}}
\makesavenoteenv{longtable}
\usepackage{graphicx}
\makeatletter
\newsavebox\pandoc@box
\newcommand*\pandocbounded[1]{% scales image to fit in text height/width
  \sbox\pandoc@box{#1}%
  \Gscale@div\@tempa{\textheight}{\dimexpr\ht\pandoc@box+\dp\pandoc@box\relax}%
  \Gscale@div\@tempb{\linewidth}{\wd\pandoc@box}%
  \ifdim\@tempb\p@<\@tempa\p@\let\@tempa\@tempb\fi% select the smaller of both
  \ifdim\@tempa\p@<\p@\scalebox{\@tempa}{\usebox\pandoc@box}%
  \else\usebox{\pandoc@box}%
  \fi%
}
% Set default figure placement to htbp
\def\fps@figure{htbp}
\makeatother



\ifLuaTeX
\usepackage[bidi=basic]{babel}
\else
\usepackage[bidi=default]{babel}
\fi
% get rid of language-specific shorthands (see #6817):
\let\LanguageShortHands\languageshorthands
\def\languageshorthands#1{}
\ifLuaTeX
  \usepackage[english]{selnolig} % disable illegal ligatures
\fi


\setlength{\emergencystretch}{3em} % prevent overfull lines

\providecommand{\tightlist}{%
  \setlength{\itemsep}{0pt}\setlength{\parskip}{0pt}}



 


\usepackage{setspace}
\onehalfspacing
\usepackage{orcidlink}
\usepackage{marvosym}
\makeatletter
\@ifpackageloaded{float}{}{\usepackage{float}}
\floatstyle{plain}
\@ifundefined{c@chapter}{\newfloat{excerpt}{h}{loexc}}{\newfloat{excerpt}{h}{loexc}[chapter]}
\floatname{excerpt}{Excerpt}
\newcommand*\listofexcerpts{\listof{excerpt}{List of Excerpts}}
\makeatother
\makeatletter
\@ifpackageloaded{tcolorbox}{}{\usepackage[skins,breakable]{tcolorbox}}
\@ifpackageloaded{fontawesome5}{}{\usepackage{fontawesome5}}
\definecolor{quarto-callout-color}{HTML}{909090}
\definecolor{quarto-callout-note-color}{HTML}{0758E5}
\definecolor{quarto-callout-important-color}{HTML}{CC1914}
\definecolor{quarto-callout-warning-color}{HTML}{EB9113}
\definecolor{quarto-callout-tip-color}{HTML}{00A047}
\definecolor{quarto-callout-caution-color}{HTML}{FC5300}
\definecolor{quarto-callout-color-frame}{HTML}{acacac}
\definecolor{quarto-callout-note-color-frame}{HTML}{4582ec}
\definecolor{quarto-callout-important-color-frame}{HTML}{d9534f}
\definecolor{quarto-callout-warning-color-frame}{HTML}{f0ad4e}
\definecolor{quarto-callout-tip-color-frame}{HTML}{02b875}
\definecolor{quarto-callout-caution-color-frame}{HTML}{fd7e14}
\makeatother
\makeatletter
\@ifpackageloaded{caption}{}{\usepackage{caption}}
\AtBeginDocument{%
\ifdefined\contentsname
  \renewcommand*\contentsname{Table of contents}
\else
  \newcommand\contentsname{Table of contents}
\fi
\ifdefined\listfigurename
  \renewcommand*\listfigurename{List of Figures}
\else
  \newcommand\listfigurename{List of Figures}
\fi
\ifdefined\listtablename
  \renewcommand*\listtablename{List of Tables}
\else
  \newcommand\listtablename{List of Tables}
\fi
\ifdefined\figurename
  \renewcommand*\figurename{Figure}
\else
  \newcommand\figurename{Figure}
\fi
\ifdefined\tablename
  \renewcommand*\tablename{Table}
\else
  \newcommand\tablename{Table}
\fi
}
\@ifpackageloaded{float}{}{\usepackage{float}}
\floatstyle{ruled}
\@ifundefined{c@chapter}{\newfloat{codelisting}{h}{lop}}{\newfloat{codelisting}{h}{lop}[chapter]}
\floatname{codelisting}{Listing}
\newcommand*\listoflistings{\listof{codelisting}{List of Listings}}
\makeatother
\makeatletter
\makeatother
\makeatletter
\@ifpackageloaded{caption}{}{\usepackage{caption}}
\@ifpackageloaded{subcaption}{}{\usepackage{subcaption}}
\makeatother
\usepackage{xcolor}
\usepackage{marginnote}
\reversemarginpar
\newcommand{\paragraphnumber}[1]{\marginnote{\color{lightgray}\tiny{#1}}[0pt]}
\usepackage{bookmark}
\IfFileExists{xurl.sty}{\usepackage{xurl}}{} % add URL line breaks if available
\urlstyle{same}
\hypersetup{
  pdftitle={Where do we draw the line?},
  pdfauthor={Zachary Batist},
  pdflang={en-CA},
  colorlinks=true,
  linkcolor={blue},
  filecolor={Maroon},
  citecolor={Blue},
  urlcolor={Blue},
  pdfcreator={LaTeX via pandoc}}


\title{Where do we draw the line?}

\author{
      {Zachary
Batist \orcidlink{0000-0003-0435-508X} \href{mailto:zachary.batist@mcgill.ca}{\Letter}} \\
          McGill University
     \\
  }

\date{February 10, 2026}
\begin{document}
\maketitle
\begin{abstract}
Researchers presently face a tension between obligations to make their
work openly accessible for public benefit, and the use of these
materials to train artificial intelligence models, which have some
questionable applications. However, this is not merely a product of AI's
arrival on the scene, and in fact represents an inherent gap between the
largely technical solutions that open science advocates propose (or
impose), and the systemic, social problems that are deeply entrenched in
research culture. This often manifests through new commitments for
participating in and maintaining the commons that are not necessarily
valued by the communities that actually make these commons possible, and
to some extent even undermine the established norms that research
communities have developed for themselves. Specifically, community
norms, which previously served to establish the boundaries that govern
who can contribute to and access the commons and in what ways, have
essentially been undermined through claims of universal access, or
claims that there should be no boundaries whatsoever.

Indeed, ``proper'' use of big data, health data, data deriving from
Indigenous sources, AI models, etc tend to be demarcated by the bounds
of established professional decorum, values and principles, as
determined by a discipline or community of practice. In this
presentation I re-cast a few recent and persistent controversies
surrounding the implementation of open information infrastructures
(e.g., the Internet Archive's National Emergency Library, use and abuse
of shadow libraries, etc) to highlight how collaborative commitments are
either maintained or re-arranged throughout distinct applications of
common resources, and why these relationships matter.
\end{abstract}


\begin{tcolorbox}[enhanced jigsaw, toptitle=1mm, titlerule=0mm, bottomtitle=1mm, colframe=quarto-callout-note-color-frame, opacitybacktitle=0.6, opacityback=0, rightrule=.15mm, colback=white, title=\textcolor{quarto-callout-note-color}{\faInfo}\hspace{0.5em}{Note}, arc=.35mm, coltitle=black, bottomrule=.15mm, left=2mm, breakable, toprule=.15mm, leftrule=.75mm, colbacktitle=quarto-callout-note-color!10!white]

This is an abstract for a paper submitted to the 13th annual gathering
of the Implementing New Knowledge Environments
(\href{https://inke.ca/}{INKE}) Partnership, whose theme this year is
``Re-Defining Open Social Scholarship in an Age of Generative
`Intelligence'\,''. The workshop will be hosted at Université de
Montréal on June 4-5, 2026, coinciding with the annual conference of the
\href{https://csdh-schn.ca/}{Canadian Society for Digital Humanities}
and the \href{https://dhsi.org/}{Digital Humanities Summer Institute}.
The full paper and slides will be accessible here in the future.\\

{This content is also available at
\href{https://zackbatist.info/inke-2026/}{zackbatist.info/inke-2026} for
a better reading experience.}

\end{tcolorbox}




\end{document}
